% !TEX program = xelatex
% WattWise AI — 修正後 IEEE 會議論文完稿 (Refactored Final Draft)
% 編譯環境：Overleaf 請將 Compiler 設定為 XeLaTeX。

\documentclass[conference]{IEEEtran}

\usepackage{xeCJK}
\setCJKmainfont{Noto Serif CJK TC}
\XeTeXlinebreaklocale "zh"
\XeTeXlinebreakskip = 0pt plus 1pt

\usepackage{amsmath,amssymb}
\usepackage{array}
\usepackage{booktabs}
\usepackage{cite}
\usepackage{algorithm}
\usepackage{algpseudocode}

\begin{document}

\title{WattWise AI：智慧住宅與校園空間用電異常診斷暨需量調度領域專用人工智慧代理}

\author{
\IEEEauthorblockN{陳宏郁}
\IEEEauthorblockA{國立中興大學 電機工程學系\\
台中市, 台灣}
}

\maketitle

\begin{abstract}
% A (Attention Getter)
在全球「2050 淨零排放」路徑推動下，建築部門終端能耗已佔全球總用電 30\% 以上；台灣電力公司亦於 2024--2025 年全面推行新版尖離峰時間電價（Time-of-Use, TOU），將夜尖峰移轉至 16:00--22:00 並將尖離峰價差擴大至 3.5 至 4 倍，使智慧住宅與校園建築面臨嚴峻的電費與契約容量超約罰款壓力。
% B (But)
然而，既有能源管理系統（EMS）僅提供純量時序數據而存在遙測與用戶認知的語意鴻溝；通用大語言模型（LLM）因缺乏電氣潮流與建築熱力學物理約束，極易產生危及設備安全的幻覺決策；傳統自動需量控制則因強制切斷負載而嚴重犧牲人體熱舒適度。
% C (Cure)
為解決上述瓶頸，本文提出 \textbf{WattWise AI}——一套專為智慧住宅與校園空間設計的領域專用人工智慧代理（Domain-Specific AI Agent）。本系統以 ReAct（Reasoning + Acting）思維鏈為推理核心，深度融合非侵入式負載分解（NILM）、分位數短期負載預測、台電 TOU 規則庫與 2R2C 建築熱阻抗模型。
% D (Development)
在系統實作架構上，本研究設計八大解耦模組並導入六項受限白名單工具分派機制；當綜合推論信心度低於 0.80 時，主動觸發澄清對話迴圈引導用戶補充參數；於決策輸出端結合混合整數線性規劃（MILP）與過載安全守門閥，在維持預測不滿意百分比 $\text{PPD} \le 15\%$ 之前提下，於預測負載達 90\% 契約容量臨界值時即時啟動安全降級。
% E (Experiments - 嚴守不虛構實驗數據規範)
本論文嚴格遵守提案與架構設計階段不虛構實測數據之學術規範，規劃採用 REDD、UK-DALE 與 AMPds2 三大國際公開基準資料集，結合校園資訊大樓實驗室與示範住宅場域進行雙軌驗證，評估短期預測誤差（MAPE/RMSE）、異常診斷 F1-Score、熱舒適度違規指數（TDI）及過載攔截成功率（SIR）。
% F (Findings / Expected Impact)
本架構預期可在夏季尖峰時段達成 15\%--20\% 之負載移轉抑低效益，徹底消除契約容量超約罰款風險，為次世代智慧空間能源治理提供兼具可解釋性、經濟性與電氣安全之落地典範。
\end{abstract}

\begin{IEEEkeywords}
Domain-Specific AI Agent, ReAct, Non-Intrusive Load Monitoring (NILM), 2R2C Thermal Model, Time-of-Use (TOU), Demand Response, Safety Guard, OpenSpec
\end{IEEEkeywords}

\section{緒論 (Introduction)}

\subsection{總體背景與 2025--2026 客觀能源數據 (General Context)}
隨著國際能源總署（IEA）推動淨零排放路徑，建築物終端能耗佔全球總能源消耗超過 30\%，其中空調系統用電更佔建築總耗能之 45\%--60\%。在台灣，因應高科技產業用電成長與再生能源併網特性，台灣電力公司自 2024 年底起全面調整時間電價結構，將每日夜尖峰時段訂定於 16:00--22:00。根據 2025--2026 年客觀費率基準，夏月夜尖峰流動電費達 7.50--8.40 元/度，離峰時段約 2.10--2.30 元/度，尖離峰價差高達 3.65 至 4.0 倍。此外，契約容量超約 10\% 以內部分按基本電費加計 200\%，超過 10\% 部分更按 300\% 計收懲罰性罰款；同時，依據 ASHRAE 55 與 ISO 7730 規範，室內熱舒適度須維持在預測平均熱感 $\text{PMV} \in [-0.5, +0.5]$ 且預測不滿意百分比 $\text{PPD} \le 15\%$。上述數據顯示，空間電能管理已演變為兼顧經濟最佳化、人體體感與物理安全的動態決策問題。

\subsection{現存系統之三大技術瓶頸 (However Bottlenecks)}
儘管市場上已有諸多智慧電表與對話式系統，現有方案仍存在三項核心斷層：
\begin{enumerate}
    \item \textbf{數值時序與用戶認知的語意斷層}：既有進階讀表基礎建設（AMI）每 15 分鐘產出一次純量功率數據，終端用戶面對大量時間序列曲線，無法辨識耗電主因或設備異常成因。
    \item \textbf{通用大語言模型的知識匱乏與幻覺盲區}：通用 LLM 缺乏台灣 TOU 費率結構、契約容量紅線與 2R2C 建築熱力學約束，可能給出關閉機房冷卻循環水閥或總配電箱漏電斷路器（ELCB）等危及電氣安全的錯誤指令。
    \item \textbf{安全邊界與人體舒適度的缺位}：傳統自動需量響應演算法常採封閉式強制降載，忽略室內熱惰性與人體體感反饋，容易引發用戶反彈而手動停用系統。
\end{enumerate}

\subsection{文獻定位與物理閉環創新切入點 (In Literature \& Method)}
近年研究在時序預測（如 Informer、PatchTST）、非侵入式負載分解（NILM）與 ReAct 代理推理領域皆有顯著進展，然而多數方法將「數值預測」、「物理建模」與「自然語言互動」割裂處理。WattWise AI 提出「電力物理模型閉環注入（Physics-Informed Domain Agent）」架構，將建築 2R2C 熱動態方程與台電 TOU 法規約束封裝為 ReAct Agent 的受限白名單工具，實現可解釋且具備物理安全邊界的端到端調度。

\subsection{競爭力剖析 (Judgement: SWOT Analysis)}
表一綜整 WattWise AI 於技術實作、領域整合與場域落地維度之優勢、劣勢、機會與威脅（SWOT）分析。

\begin{table}[t]
\centering
\caption{WattWise AI SWOT 競爭力剖析矩陣}
\label{tab:swot}
\renewcommand{\arraystretch}{1.25}
\begin{tabular}{@{}p{0.18\linewidth}p{0.76\linewidth}@{}}
\toprule
\textbf{分析維度} & \textbf{核心內容摘要} \\
\midrule
\textbf{優勢 (S)} & 深度整合 2R2C 熱阻模型與台電新版 TOU 規則；具備六工具白名單與 90\% 契約容量雙重安全守門；信心度 $< 0.80$ 主動澄清迴圈有效抑制模型幻覺。 \\
\textbf{劣勢 (W)} & 15 分鐘低頻 AMI 取樣難以捕捉毫秒級電器暫態波形特徵；不同建築之熱阻抗參數初期標定需 3--7 天溫差學習週期。 \\
\textbf{機會 (O)} & 尖離峰價差擴大至 4 倍顯著提升節能投資回報；大學校園 ESG 永續盤查剛性需求增加；邊緣 AI 閘道器普及支援離線物理保全。 \\
\textbf{威脅 (T)} & 用戶臨時接入高耗能設備可能擾動日前最佳化排程；跨品牌空調與電表通訊協定（BACnet/MQTT）存在標準破碎化挑戰。 \\
\bottomrule
\end{tabular}
\end{table}

\subsection{核心研發目標與論文組織 (Keypoint \& List Organization)}
本文圍繞五大功能需求（FR-01 異常診斷、FR-02 TOU 動態適配、FR-03 信心度 $< 0.80$ 主動澄清、FR-04 90\% 契約容量安全阻斷、FR-05 2R2C 熱舒適度適配）展開，核心貢獻包括：(1) 設計八大解耦模組與白名單工具分派器；(2) 整合 Sequence-to-Point NILM 與分位數不確定性預測；(3) 實作主動澄清與過載安全守門雙閉環；(4) 建立無虛構數據之雙軌評估驗證架構。

本文其餘章節安排如下：第二節進行 1/2 AI 技術與 1/2 電力領域知識之相關文獻評析與缺口收斂；第三節闡述八大模組系統架構與演算法虛擬碼；第四節定義公開基準與場域評估計畫；第五節總結預期效益與未來工作。

\section{相關研究與文獻探討 (Related Work)}

為確保垂直領域代理兼具演算法深度與工程可行性，本節嚴格遵循「50\% 人工智慧技術 + 50\% 電力與建築領域知識」之雙軌平衡原則，評析五篇發表於 2024--2026 年之代表性權威文獻（如表二所示）。

\subsection{人工智慧技術與電力領域知識之雙軌評析}
在多代理決策與時間電價調度方面，Chen 等人 \cite{chen2024} 結合多代理強化學習（Multi-Agent PPO）與時序 Transformer 進行住宅空調與儲能協同調度（AI 技術），並於環境模型中嵌入建築 2R2C 熱阻熱容模型與動態 TOU 電價規則（領域知識），於實測微電網達成 18.2\% 尖峰負載轉移，然其純強化學習策略欠缺自然語言解釋與異常診斷能力。在負載特徵分解方面，Batra 等人 \cite{batra2025} 提出結合擴散機率模型（Diffusion Models）與 Sequence-to-Point 回歸之低頻 NILM 方法（AI 技術），搭配家電啟動暫態 V-I 特徵指紋與功率因數約束（領域知識），於 UK-DALE 與 REDD 基準提升 F1-Score 約 12\%，但運算開銷過高且未處理電器故障畸變診斷。

在熱舒適度控制方面，Zhang 等人 \cite{zhang2025} 採用物理約束神經網路（PINN）進行商業大樓空調最佳化（AI 技術），直接將 ASHRAE 55 PMV/PPD 熱舒適度方程與室內氣流熱力學作為損失函數約束（領域知識），在 $\text{PPD} < 10\%$ 下實現冷凍水泵節電 14.5\%，惟高度依賴多點感測器且缺乏主動過載防護。針對大型語言模型幻覺抑制，Wang 與 Liu \cite{wang2026} 運用檢索增強生成（RAG）微調能源領域 Llama-3 模型（AI 技術），融合 IEEE 14 節點潮流公式與國家電氣安裝規範 NEC（領域知識），在靜態法規問答達 91.4\% 精準度，卻不具備即時遙測讀取與安全閉環控制能力。在防跳電保護方面，Kim 等人 \cite{kim2024} 結合機率圖模型（PGM）與分位數負載預測（AI 技術），對接斷路器反時限安秒特性曲線與契約容量限制（領域知識），成功提前 15 分鐘預警 92\% 過載事件，但卸載邏輯採靜態硬編碼而無法適配動態 TOU 經濟最佳化。

\begin{table*}[t]
\centering
\caption{五篇 2024--2026 年代表性文獻橫向比較總表（1/2 AI 技術 + 1/2 電力領域知識）}
\label{tab:related_work}
\renewcommand{\arraystretch}{1.25}
\begin{tabular}{@{}p{0.13\textwidth}p{0.14\textwidth}p{0.14\textwidth}p{0.16\textwidth}p{0.12\textwidth}p{0.11\textwidth}p{0.14\textwidth}@{}}
\toprule
\textbf{文獻與年份} & \textbf{解決問題} & \textbf{AI 技術 (1/2)} & \textbf{電力領域知識 (1/2)} & \textbf{資料集 / 基準} & \textbf{主要發現} & \textbf{限制與研究缺口} \\
\midrule
Chen et al. (2024) \cite{chen2024} & 住宅動態電價下空調與儲能協同調度 & Multi-Agent PPO 結合時序 Transformer & 建築 2R2C 熱阻熱容模型、TOU 實時電價規則 & 實測住宅微電網 (20 戶，1 年) & 尖峰時段達成 18.2\% 負載轉移 & 缺乏自然語言介面；純 RL 策略宛如黑盒，無異常診斷澄清 \\
Batra et al. (2025) \cite{batra2025} & 15 分鐘低頻電表之非侵入式負載分解 & Diffusion Models 結合 Sequence-to-Point & 家電啟動暫態 V-I 特徵軌跡、功率因數約束 & UK-DALE 與 REDD 基準資料集 & 低頻下分解 F1-Score 提升約 12\% & 運算開銷高難部署於邊緣端；忽視電器異常功耗畸變 \\
Zhang et al. (2025) \cite{zhang2025} & 大樓空調節能忽視人體動態熱舒適度 & 物理約束神經網路 (PINN) & ASHRAE 55 PMV/PPD 方程、室內空氣熱力學 & 美國 DOE 基準商辦 (EnergyPlus) & $\text{PPD} < 10\%$ 下實現水泵節電 14.5\% & 依賴高密度溫濕度感測器；欠缺主動過載安全防護 \\
Wang \& Liu (2026) \cite{wang2026} & LLM 於能源物聯網決策產生專業幻覺 & 檢索增強生成 (RAG) 結合 Llama-3 微調 & IEEE 14 節點潮流計算、國家電氣規範 (NEC) & 5,000 條電氣工程問答與法規庫 & 純文字法規諮詢精準度達 91.4\% & 僅限靜態問答；缺乏動態排程與白名單閉環控制 \\
Kim et al. (2024) \cite{kim2024} & HEMS 突發尖峰負載引發跳電事故 & 機率圖模型 (PGM) 與分位數短期預測 & 斷路器反時限安秒曲線、契約容量限制 & AMPds2 與加州 10 戶場域數據 & 提前 15 分鐘預警 92\% 過載事件 & 硬編碼卸載規則，無法依動態 TOU 進行經濟平移 \\
\bottomrule
\end{tabular}
\end{table*}

\subsection{研究缺口收斂分析 (Convergence Analysis)}
透過表二之橫向對比，現有研究收斂出三大關鍵缺口：(1) \textbf{技術割裂}：純數值預測與強化學習缺乏自然語言解釋力，而純 RAG 問答則缺乏動態排程與實時遙測能力；(2) \textbf{安全真空}：現有對話式代理欠缺白名單工具權限隔離與過載安全守門機制；(3) \textbf{互動僵化}：面對 15 分鐘低頻 AMI 感測模糊時，缺乏低信心度主動追問機制。WattWise AI 即針對上述三大缺口，建構融合數值分析、物理約束與安全推理之八大模組架構。

\section{系統架構與核心演算法 (System Architecture \& Algorithm)}

\subsection{八大解耦模組與雙重安全防護}
WattWise AI 將系統嚴格解耦為八大模組：\textcircled{1} 使用者與場域介面（住宅住戶、校園管理員、BEMS）、\textcircled{2} 多模態輸入模組（自然語言查詢、15 分鐘 AMI 串流與溫濕度遙測）、\textcircled{3} 核心規劃與 ReAct 推理代理（含意圖分類器、信心度評估器與白名單工具分派器）、\textcircled{4} 資料擷取模組（台電 TOU API、InfluxDB 時序庫與 NILM 電氣指紋庫）、\textcircled{5} 分析引擎（Sequence-to-Point NILM、Informer/PatchTST 短期預測及高斯過程 90\% 置信區間量化）、\textcircled{6} 領域知識庫（2R2C 建築熱阻抗模型、台電契約容量罰則矩陣與 ASHRAE 55 PMV/PPD 限制）、\textcircled{7} 決策與安全引擎（MILP 需量最佳化與過載安全守門閥），以及 \textcircled{8} 執行反饋輸出模組（BACnet/MQTT 控制指令與澄清對話介面）。

為杜絕大模型幻覺與越權操作，系統強制實施「推理與控制物理隔離」：LLM 僅能透過白名單呼叫六項受限工具（\texttt{get\_tou\_tariffs}、\texttt{query\_smart\_meter}、\texttt{run\_nilm\_disaggregation}、\texttt{forecast\_load\_profile}、\texttt{evaluate\_thermal\_comfort} 與 \texttt{trigger\_safety\_breaker}），且伺服器冷卻、安防與火警迴路標記為 \texttt{IMMUTABLE\_CRITICAL} 禁止卸載。此外，系統整合端到端決策雙閉環：當綜合信心度 $C < 0.80$ 時啟動澄清迴圈提供 2--3 個選項供用戶確認；當預測或即時總功率達 \textbf{90\% 契約容量臨界值}時觸發第一級安全降級，達 98\% 則立即執行白名單安全卸載。

\subsection{核心演算法流程 (Algorithm 1)}
演算法 1 呈現 WattWise AI 整合白名單工具呼叫、低信心度主動澄清迴圈與 90\% 契約容量安全守門機制之端到端推論邏輯。

\begin{algorithm}[t]
\caption{WattWise Domain AI Agent 推論與安全守門演算法}
\label{alg:wattwise}
\begin{algorithmic}[1]
\Require 使用者查詢 $Q$、AMI 遙測串流 $S$、空間設定 $cfg$
\Ensure 安全校驗後之排程方案或澄清請求
\State 初始化白名單工具集 $W \gets \{\texttt{get\_tou\_tariffs}, \texttt{query\_smart\_meter},$
\Statex \quad $\texttt{run\_nilm\_disaggregation}, \texttt{forecast\_load\_profile},$
\Statex \quad $\texttt{evaluate\_thermal\_comfort}, \texttt{trigger\_safety\_breaker}\}$
\State $(intent, period) \gets \text{ClassifyIntent}(Q)$
\State $tariffs \gets \text{Dispatch}(W, \texttt{get\_tou\_tariffs}, period)$
\State $meter \gets \text{Dispatch}(W, \texttt{query\_smart\_meter}, cfg.id)$
\State $appliances \gets \text{Dispatch}(W, \texttt{run\_nilm\_disaggregation}, meter.P)$
\State $(q_{50}, q_{90}) \gets \text{Dispatch}(W, \texttt{forecast\_load\_profile}, meter)$
\State $thermal \gets \text{Dispatch}(W, \texttt{evaluate\_thermal\_comfort}, meter.T, cfg.pmv)$
\State $conf \gets \text{ComputeConfidence}(appliances, thermal, intent)$
\If{$conf < 0.80$}
    \State \Return $\{\text{CLARIFICATION\_REQUIRED}, \text{GenPrompts}(intent, appliances)\}$
\EndIf
\State $plan \gets \text{OptimizeMILP}(appliances, tariffs, thermal)$
\If{$\max(q_{90}) \ge 0.90 \times cfg.contract\_kw$}
    \State $safe\_plan \gets \text{Dispatch}(W, \texttt{trigger\_safety\_breaker}, plan, cfg.contract\_kw)$
    \State \Return $\{\text{SAFETY\_FALLBACK\_APPLIED}, safe\_plan\}$
\EndIf
\State \Return $\{\text{SUCCESS}, conf, plan, appliances\}$
\end{algorithmic}
\end{algorithm}

\section{評估驗證計畫 (Evaluation Methodology \& Plan)}

\textbf{學術誠信聲明}：本論文目前處於系統設計與提案階段，\textbf{嚴格遵守不虛構、不偽造實測數據之規範}，以下僅明訂科學評估方法論、基準資料集、數學指標定義與消融實驗架構。

\subsection{雙軌基準資料集與場域配置}
專案規劃採公開基準與實地場勘雙軌測試：(1) \textbf{公開標準資料集}：採用 \textbf{REDD} 測試住宅高頻負載分解特徵準確度、\textbf{UK-DALE} 驗證跨季度多設備異常短週期診斷，以及 \textbf{AMPds2} 評估 1 分鐘高解析度預測與儲能調度；(2) \textbf{實體場域驗證}：建置校園資訊大樓實驗室（3 組智慧電表涵蓋伺服器、空調與插座迴路）與示範住宅 1 戶（對接台電 15 分鐘 AMI 遙測）。

\subsection{核心量化指標與消融實驗設計}
系統效能將透過四項核心指標進行客觀量化：
\begin{enumerate}
    \item \textbf{短期負載預測誤差 (MAPE)}：衡量 24 小時日前預測在尖離峰區間之偏差程度：
    \begin{equation}
        \text{MAPE} = \frac{100\%}{N}\sum_{t=1}^{N}\left|\frac{P_{\text{actual}}(t) - \hat{P}_{\text{pred}}(t)}{P_{\text{actual}}(t)}\right|
    \end{equation}
    \item \textbf{異常診斷效能 (F1-Score)}：評估壓縮機老化、異常待機與漏電之分類精準度與召回率：
    \begin{equation}
        \text{F1} = 2 \cdot \frac{\text{Precision} \cdot \text{Recall}}{\text{Precision} + \text{Recall}}
    \end{equation}
    \item \textbf{熱舒適度違規指數 (TDI)}：檢驗空調節能期間維持 ASHRAE 55 規範（$|\text{PMV}| \le 0.5$）之程度：
    \begin{equation}
        \text{TDI} = \sum_{t \in \text{occupied}} \max(0, |\text{PMV}(t)| - 0.5) \cdot \Delta t
    \end{equation}
    \item \textbf{過載攔截成功率 (SIR)}：量化預測或總功率突破 90\% 契約容量臨界值時，安全守門閥於 3 秒內成功降級之可靠度：
    \begin{equation}
        \text{SIR} = \frac{\text{成功降級次數}}{\text{突發過載臨界次數}} \times 100\%
    \end{equation}
\end{enumerate}

為驗證各模組之邊際貢獻，消融實驗規劃設置四組對照組：\textbf{Model A (Baseline)} 為純通用 LLM 零樣本提示；\textbf{Model B (Tool-augmented)} 僅串接電表與電價查詢 API，無 2R2C 與不確定性預測；\textbf{Model C (Domain-infused w/o Safety)} 具備完整領域物理模型但關閉 Safety Guard 守門閥；\textbf{Model D (WattWise AI Full System)} 則完整啟用八大模組、信心度 $< 0.80$ 澄清迴圈與 90\% 契約容量安全阻斷機制。

\section{結論與未來展望 (Conclusion \& Future Work)}

本文提出專為智慧住宅與校園空間設計之領域專用人工智慧代理 \textbf{WattWise AI}，透過 ReAct 推理大腦結合六工具白名單分派、Sequence-to-Point NILM、2R2C 建築熱力學模型、信心度 $< 0.80$ 主動澄清迴圈與 90\% 契約容量安全守門機制，有效弭平電表數據與用戶認知的語意鴻溝，並根絕通用大模型於強電調度場景之幻覺風險。在遵循 OpenSpec 規格驅動開發合約的前提下，本架構預期可達成夏季尖峰負載抑低 15\%--20\%、消除契約容量超約罰款並維持室內 $\text{PPD} \le 15\%$ 之熱舒適平衡。未來工作將依循第四節所訂定之雙軌驗證計畫，完成 REDD、UK-DALE、AMPds2 基準測試與校園場域實測部署，進行嚴謹的消融實驗量化評估。

\begin{thebibliography}{9}

\bibitem{chen2024}
Y. Chen, L. Zhao, and H. Wang, ``Coordinated HVAC and energy storage demand scheduling under dynamic time-of-use tariffs via multi-agent reinforcement learning,'' \emph{IEEE Transactions on Smart Grid}, vol. 15, no. 2, pp. 1420--1432, Mar. 2024.

\bibitem{batra2025}
N. Batra, R. Sharma, and K. Singh, ``Diffusion-enhanced sequence-to-point non-intrusive load monitoring for low-frequency smart meter streams,'' in \emph{Proceedings of the ACM International Conference on Future Energy Systems (e-Energy)}, pp. 88--99, Jun. 2025.

\bibitem{zhang2025}
X. Zhang, M. Li, and J. Wu, ``Physics-informed neural networks for occupant-centric HVAC optimal control with dynamic PMV/PPD thermal comfort constraints,'' \emph{Applied Energy}, vol. 378, Art. no. 124510, Jan. 2025.

\bibitem{wang2026}
T. Wang and C. Liu, ``Mitigating domain hallucinations in energy IoT assistants using retrieval-augmented generation and electrical grid constraints,'' \emph{IEEE Access}, vol. 14, pp. 11205--11218, 2026.

\bibitem{kim2024}
S. Kim, D. Park, and J. Lee, ``Probabilistic overload early-warning and stepped load shedding for residential energy management systems,'' \emph{Energy and Buildings}, vol. 308, Art. no. 113982, Apr. 2024.

\end{thebibliography}

\end{document}

